\scanpage{89}
\mspage{38}

pour vérifier de telles propriétés, c'est soit par la $W$-propreté de $f$, soit par la $W$-lissité -- ce qui revient à vérifier que certaines catégories bien précisées sont $W$-asphériques, sans avoir à se soucier de chgts de base et de vérifier que des flèches construites en termes de chgt. de base arbitraires sont dans $W$. Et l'emploi de la $W$-propreté et de la $W$-lissité est d'autant plus sympa, que ces notions sont stables par composition, alors que \deleted{\unsure{toutes}} les notions \unsure{dérivées}, ainsi que les notions de fibration à la clef, ont une fâcheuse tendance à ne \uline{pas} être stables par composition.

On pourrait appeler pseudo-\inserted{$W$-}propres resp. quasi-\inserted{$W$-}propres (on \unsure{verra} mieux que pseudo) ce qu'on vient de dégager) -- j'entends des \unsure{ricanements}! Et si on veut les deux à la fois -- $\lambda$ et $\lambda^\circ$ \unsure{associés} comme on a dit? Je vois le $\unsure{=}$ principe de description avec \inserted{des} $S''\to S'$ sur $S$, d'exiger que $X''\mathbin{\unsure{\text{---}|\to}}X'$ soit transformé en des $\in W$.

\unsure{Avec|Donc|D'où} les choix suivants qui ont l'air sympa.

% archive sheet number: 89
